% ===========================================================================
% DISCUSSION: revision of 5 Oct 2026 (Aron; in Overleaf since 6 Oct 2026).
% Changes from the 4 Oct Overleaf version:
%   * 5.1 para 2: frontier defector runs now reported (agents send defectors more with
%     myths), so they move from "would test this" to evidence.
%   * 5.1 para 3 rewritten: the myth as the writer's plan, tied to promises and
%     implementation intentions; scope of the cultural reading stated.
%   * 5.1 para 4 rewritten: the monitoring argument no longer claims the influence is
%     invisible in the text (the replay shows the stated amount carries it).
%   * 5.2 adds: replay is single decisions; plan-vs-myth untested; one reader per myth;
%     pilots n = 3; Holm; n = 10 extension ran with both groups together.
%   * 5.3 updated to the new framing.
%   * Clarity pass (5 Oct, evening): 5.1 reordered and split; writing-first claim scoped to
%     single-model groups (matches Results 4.4); US spelling.
% Cut pass (5 Oct, late): noise speculation cut to one sentence; plan paragraph no longer
%   restates the replay numbers; Limitations tightened (830 -> 756 prose words).
% 6 Oct afternoon (with Aron): style pass (fewer colons/semicolons and metaphors).
% 8 Oct: second editorial review N060-N074.
% 8 Oct (Ed, via Aron): replay mentions removed (5.1, safety paragraph, Limitations, Conclusion).
% 8 Oct: editorial review E113-E126 (E122 maxim cut; limitations split in two).
% ===========================================================================
% ivar: 
% arabella "transmission modulates cooperation because; some form of linguistic transmission occurs, agents adapt language use, and myths increase cooperation. so, over time, the agents' own myths also change as a result of their partners language (?). " 
% transmission is highest when there is conflict between agents', one agents' myth is cooperative, the other one not. 
% also, recent claude's myths have more transmission because they take up rules of other myths they don't agree with in order to refute them. 

Overall, our experiments reveal patterns of agent behaviour consistent with cultural transmission in humans: agent populations that share myths exhibit greater cooperation than those that do not. We also observe that this effect is amplified in mixed groups, where less cooperative agents can be positively influenced by other, more cooperative agent behavior. 
Interestingly, while the language of an agent is most influential when it comes to its own game-play behaviour, agents do pick up and reuse the language of their peers, thus, over time shared myths evolve to be artifacts jointly created by the agent population. Taken together, these results suggest that shared language, and its transmission and evolution, can play a similarly important role in populations of agents as it does for humans when it comes to shaping cooperation. 

\subsection{5.1 Interpretation}
\label{discussion-interpretation}

Behaviorally, our agents resemble the conditional cooperators of human public-goods experiments, who give more to partners who give more~\citep{fischbacher2001}. GPT-5 Nano plays the free rider, and its partners stop sending within a few rounds. One possibility, which we have not tested, is that a model that under-corrects for noise it was warned about drifts downward, as imperfect conditional cooperators do~\citep{fischbacher2010}. Tolerating apparent shortfalls, as GPT-5.6 Sol appears to, helps sustain cooperation under noise~\citep{wu1995}. % TODO(Results): add the noise-correction analysis with uncertainty. Game-only dyads: Sonnet 4.5 sends $0.44 more than it saw its partner send, Opus 5 $0.59 more; compare with the realized shortfall in those runs (expected $0.50; the 210 distinct noise draws average $0.51).

The cooperative models' lower cooperation with GPT-5 Nano shows up in what they send; their mean return proportions stay similar across partners (Section~4.2). In human groups, costly punishment of free riders sustains cooperation that conditional cooperation alone lets decay~\citep{fehr2000,fehr2002,boyd2003}. Our game has no separate costly punishment, so as a sender an agent can limit its exposure to a withholder only by sending less, which may be how a withholding model lowers cooperation in mixed groups. With myths, frontier non-defectors also sent scripted defectors more (supplementary material). Whether LLM agents would sanction defectors when a game allows it is a question for designs with an explicit punishment option. % [Aron: our Aug. punishment pilots (researchlog 21--23 Aug) found Gemini deducts selectively when given an explicit option, hence "spontaneously"]

That agents' opening transfers track the amounts their founding myths name (Section~4.4) has a human parallel. A promise made before a trust game raises the promiser's own trustworthiness~\citep{charness2006}, and a concrete plan raises follow-through~\citep{gollwitzer1999}. It also fits how writing first helps single-model groups, where the amount a founding myth names tracks the opening transfer and reciprocity may carry it forward. In mixed groups, adjusting for the opening does not shrink the advantage of writing first. The same result limits a cultural reading of the treatment. In our design, we find clearer evidence that myths pass words between agents than that they pass on amounts or morals readers act on. Agents' own myths track their transfers, but the bundled treatment does not tell us whether writing, reading or the accompanying instruction produces the gain. Whether a plain written plan would work as well as a myth remains open.

Two results matter for multi-agent safety. First, a simple writing task raised cooperation, especially where a withholding model was present. Second, what tracked behavior was concrete and visible in the text: the amount an agent's founding myth names went with its opening transfer, whereas the judge's moral labels predicted little of later play. This motivates testing monitors that read for explicit amounts and plans rather than themes, including on channels that agents improvise for themselves, like the message board in the Hugging Face incident; our study does not show how well such monitors would work. Whether influence between agents can also run through content such a monitor would not flag~\citep{motwani2024} is an open question.

\subsection{5.2 Limitations}
\label{limitations}

The myth condition combines writing, reading a partner's myth, an instruction to use myths, and a longer context. We did not separate these components over whole runs or compare myths with plain written plans. Each myth has one reader, who builds on its own last myth, and agents keep only their last three myth exchanges in context (their own myths and the partner myths quoted in the prompts), so our weak evidence for behavioral transmission applies to this channel. The myth prompt also asks agents to use their own previous myth as inspiration, which may favor their own text over a partner's. The design does not separate authorship from where a text sits in the conversation or from the instruction to build on one's own previous myth.

Conclusions about adding a withholding model rest on GPT-5 Nano alone, and population contrasts with five to ten runs per cell are exploratory. The main comparisons use one game, one prompt set and one noise level, and comparisons between models also compare their request settings. LLM judges labeled myth morals with moderate agreement ($\kappa=0.53$) and without human validation, and we measure cooperation only within the game.

\subsection{5.3 Conclusion}
\label{conclusion}

In our trust game, both the myth task and the mix of models affect how much LLM agents cooperate. The myth task raises cooperation, most of all in groups that contain a withholding model, which otherwise lowers its partners' cooperation. Agents' transfers track the amounts written into their own myths far more than those in myths they read, and we find clearer evidence that words pass between agents than that behavior does. Agents should therefore be evaluated with the kinds of partners and channels they will meet, including channels where their own earlier text stays in context.
